% ID: SEL-D0033-OB-E04
% ESTADO: cotejada. Segunda pasada visual y borrador compilado revisados.
% Fuente: Selectividad/matematicas_aplicadas/titular_septiembre/Examen-Matemáticas Aplicadas a las Ciencias Sociales.pdf
% SHA-256: 0f95b6aefa78f1a13e2115b281338bd807177671c303c380180957fe03f0a605
% Página: 2; opción B; ejercicio original 4.

La edad de los empleados de una empresa sigue una ley Normal de varianza 64 y media desconocida. Se toma una muestra aleatoria simple de 16 empleados de dicha empresa, obteniéndose las siguientes edades
\[
\begin{gathered}
30\quad42\quad38\quad45\quad52\quad60\quad21\quad26\\
33\quad44\quad28\quad49\quad37\quad41\quad38\quad40
\end{gathered}
\]
\begin{enumerate}[label=\alph*)]
\item \textbf{(1.5 puntos)} Obtenga un intervalo de confianza para estimar la edad media de los empleados, con un nivel de confianza del 97\%.
\item \textbf{(1 punto)} Calcule el tamaño mínimo de la muestra que se ha de tomar para estimar la edad media de los empleados, con un error inferior a 2 años y un nivel de confianza del 99\%.
\end{enumerate}
